\subsection{复测度的像}

设 \(\theta\) 是 T 上的一个复测度，并设 \(\pi\) 是一个由 T 到局部紧空间 X 中的映射；设 \(\pi\) 是 \(\theta\) -可测的，且对每个 \(f \in \mathcal{K}(X; \mathbf{C})\) ，\(f \circ \pi\) 是本质 \(\theta\) -可积的。由于说一个函数关于 \(\theta\) 可测（相应地，本质可积）与说它关于 \(|\theta|\) 可测（相应地，本质可积）是一回事，这就表明 \(\pi\) 是 \(|\theta|\) -逆紧的。若 \(f \in \mathcal{K}(X; \mathbf{C})\) ，

\[
\left| \int (f \circ \pi) d \theta \right| \leqslant \int (| f | \circ \pi) d | \theta |;\tag{7}
\]

立即推出，线性形式 \(f \mapsto \int(f \circ \pi) d\theta\) （\(\mathcal{K}(X; \mathbf{C})\) 上的）是 X 上的一个复测度（Ch. III, §1, No.~3, 命题 6），于是可以给出下列定义：

定义 2. --- 设 \(\theta\) 是局部紧空间 T 上的一个复测度。称一个由 T 到局部紧空间 X 中的映射 \(\pi\) 为 \(\theta\) -逆紧的，如果它是 \(|\theta|\) -逆紧的。测度 \(f \mapsto \int(f \circ \pi) d\theta\) 这时称为 \(\theta\) 在 \(\pi\) 下的像，记作 \(\pi(\theta)\) 。

于是关系式 (7) 可以写成下列形式：

\[
| \pi (\theta) | \leqslant \pi (| \theta |).\tag{8}
\]

测度 \(\pi(\theta)\) 可以是零而 \(\theta\) 不是零，只要取 T 为缩成两点 a, b 的空间，取 \(\theta\) 为测度 \(\varepsilon_{a} - \varepsilon_{b}\) ，并取 \(\pi\) 为一个常值映射，即可立即看出这一点。

设 \(\theta\) 和 \(\theta'\) 是 T 上的两个复测度；若 \(\pi\) 是 \(\theta\) -逆紧的且是 \(\theta'\) -逆紧的，则由命题 6 的推论 2 推出 \(\pi\) 是 \((\theta + \theta')\) -逆紧的，因为 \(|\theta + \theta'| \leqslant |\theta| + |\theta'|\) ，且显然 \(\pi(\theta + \theta') = \pi(\theta) + \pi(\theta')\) 。

特别地，若 \(\theta\) 是一个实测度且 \(\pi\) 是 \(\theta\) -逆紧的，则

\[
\pi (\theta) = \pi (\theta^ {+}) - \pi (\theta^ {-}).\tag{9}
\]

前面建立的若干结果可以立即推广到复测度；我们列举其中最重要的。

命题 7. --- 设 \(\theta\) 是 T 上的一个复测度，\(\pi\) 是一个由 T 到局部紧空间 X 中的 \(\theta\) -逆紧映射，\(\nu\) 是像测度 \(\pi(\theta)\) 。

\begin{enumerate}
\def\labelenumi{\alph{enumi})}
\item
  设 A 是 X 的一个子集；若 \(\overline{\pi}^{1}(A)\) 是局部 \(\theta\) -可忽略的，则 A 是局部 \(\nu\) -可忽略的。
\item
  设 \(f\) 是一个由 \(X\) 到拓扑空间中的映射；若 \(f \circ \pi\) 是 \(\theta\) -可测的，则 \(f\) 是 \(\nu\) -可测的。
\item
  设 \(\mathbf{f}\) 是定义在 \(X\) 上、取值于一个 Banach 空间 \(F\) 中的函数；若 \(\mathbf{f} \circ \pi\) 是本质 \(\theta\) -可积的，则 \(\mathbf{f}\) 是本质 \(\nu\) -可积的，且
\end{enumerate}

\[
\int \mathbf {f} (\pi (t)) d \theta (t) = \int \mathbf {f} (x) d \nu (x).\tag{10}
\]

注意到公式 (8)，这些结果可由命题 2 的推论 2、命题 3 以及 No.~2 的定理 1 得出。

